\section{Introduction}

As large language models (LLMs) move from pre-training to deployment, a new bottleneck emerges: a single base model now produces many post-training variants---quantized (GPTQ~\cite{frantar2022gptq}, GGUF, BitsAndBytes~\cite{dettmers2022gpt3}), LoRA-adapted~\cite{hu2022lora}, or distilled---that must be evaluated before release~\cite{grattafiori2024llama, qwen2025qwen3}. Existing evaluation largely relies on aggregate accuracy or perplexity, which can reveal that a variant has degraded but not \emph{why}. As a result, developers often resort to costly trial-and-error when debugging low-bit quantization failures, catastrophic forgetting, or prediction-head corruption. What is missing is a diagnostic that not only predicts degradation, but also identifies which component of the model has drifted from the base checkpoint.

A natural approach is to compare internal representations. Prior representational similarity methods such as SVCCA~\cite{raghu2017svcca}, CKA~\cite{kornblith2019similarity}, and generalized shape metrics summarize two activation spaces with a single alignment score. However, these metrics are descriptive rather than predictive: they are not tied to downstream cross-entropy risk and provide little insight into the mechanism of failure~\cite{ding2021grounding, davari2023reliability, klabunde2023similarity}. For deployment-time diagnostics, we instead need a representation-level quantity that both bounds downstream risk and decomposes into interpretable failure modes.

In this work, we show that modern LLMs admit exactly such a decomposition. Our analysis builds on two structural properties of Transformer-based LLMs. First, LLMs naturally separate into a non-linear backbone and a linear prediction head (\texttt{lm\_head}), allowing feature-side and head-side errors to be analyzed independently. Second, prior work suggests that hidden representations across related LLMs are approximately linearly aligned and near-isometric~\cite{park2023linear, moschella2022relative}. Leveraging these properties, we introduce \textbf{PRISM} (\textbf{P}roxy \textbf{R}isk \textbf{I}nference via \textbf{S}tructural \textbf{M}apping; Fig.~\ref{fig:prism_concept}), a closed-form upper bound on the cross-entropy risk gap between a target model and its variant.

PRISM decomposes model drift into three independently measurable axes: \emph{scale mismatch} $\Delta \rho$, which captures activation magnitude collapse; \emph{shape mismatch} $1-\Omega$, which captures distortion of the feature geometry; and a covariance-weighted \emph{head discrepancy} $\gamma$, which measures prediction-head divergence. Unlike scalar similarity scores, these axes correspond to distinct empirical failure modes. For example, low-bit quantization primarily induces shape distortion, whereas quantizing the output projection directly inflates head divergence. We further derive a sharper feature-side Lipschitz constant for cross-entropy that remains informative even at the LLM vocabulary scale (Sec~\ref{subsec:unified_bound}), enabling the resulting bound to track downstream degradation in practice.

Beyond post-hoc diagnosis, PRISM also provides a training signal. Under frozen-head LoRA fine-tuning, the head discrepancy term vanishes, leaving the differentiable shape term as a clean regularization target for backbone drift. This allows us to regularize fine-tuning directly by penalizing feature-geometry distortion, reducing catastrophic forgetting without replay (Sec.~\ref{subsec:shape_reg_exp}).
Empirically, across Llama-, Qwen-, Ministral-, and DeepSeek-based variants and five benchmarks, PRISM consistently tracks degradation across both post-training quantization and LoRA fine-tuning. The bound achieves mean Spearman correlations of $0.820$ for quantized variants and $0.831$ for LoRA forgetting, while its decomposition localizes distinct failure mechanisms across quantization protocols and fine-tuning tasks. Moreover, the proposed shape regularizer outperforms replay-based baselines in mitigating forgetting.

Our contributions are three-fold.
\begin{enumerate}[leftmargin=*,nosep,topsep=2pt]
    \item \textbf{Theory: a closed-form CE risk bound with three diagnostic axes.} $|\mathcal{R}_T - \mathcal{R}_P|$ admits a closed-form upper bound as a sum of three axes: scale $(\Delta\rho)^2$ and shape $2\rho_T\rho_P(1{-}\Omega)$ (from an exact Procrustes residual decomposition), plus a covariance-weighted head-discrepancy term (Theorem~\ref{thm:unified_bound}; Fig.~\ref{fig:prism_concept}).

    \item \textbf{Framework: a unified diagnostic that doubles as a training objective.} Computed from features and head weights alone, the bound applies across both PTQ and frozen-head LoRA; the differentiable trace form $\Omega$ further turns it into a training-time regularizer against catastrophic forgetting (Sec.~\ref{subsec:action}).

    \item \textbf{Empirical: rank-consistent across both settings, with axis-level localization.} On two model families (Ministral and DeepSeek in Appendix~\ref{app:per_model_tables}) over five benchmarks, the bound ranks PTQ variants ($r_s{=}0.820$) and LoRA checkpoints ($r_s{=}0.831$) consistently (\emph{predictiveness}); its axes separate failure modes---shape distortion at Q2/Q3 PTQ, head divergence at Qwen3 Q6\_K/Q8\_0, scale-axis separability under cross-task LoRA drift (\emph{decomposability}); and a shape regularizer outperforms experience replay at mitigating downstream forgetting (\emph{actionability}).
\end{enumerate}
